\chapter{Long Context and RAG Systems}
\label{chap:long_context_rag}

\begin{tldr}
Context windows grew from 4K to 128K--1M while corpora stayed larger than any of them, and this chapter is the program's canonical home for the two questions that follow. \textbf{How a window is actually built and verified}: continued pretraining with progressive length extension (the frontier recipe), RoPE-scaling retrofits---position interpolation, NTK-aware scaling, YaRN---as the no-retrain path, and the claimed-versus-effective-length evaluation that decides whether the advertised window is real. \textbf{When to spend that window instead of retrieving}: the context-versus-retrieval decision, its cost arithmetic, and the hybrid that most production systems land on. The RAG \emph{pipeline} is not this chapter's material---it is canonical in [NLP~9] (chunking, retrieval-paradigm choice, query transformation, grounded generation, advanced patterns, RAG evaluation, failure taxonomy), with the retrieval infrastructure under it in Volume~III: vector-index internals [SR~3], bi-encoders, rerankers and fusion [SR~4], and production retrieval serving [SR~8]. What is kept here of RAG is exactly what the architectural decision needs---what a pipeline costs to own and where it breaks---so the trade can be made with numbers instead of taste.
\end{tldr}

%==============================================================================
\section{Long Context Techniques}
\label{sec:long_context}
%==============================================================================

\subsection{The Problem: Extending Trained Context Length}

LLMs are trained with a fixed maximum sequence length---4K--8K for the 2023 generation of open models, and even for today's frontier models the bulk of pretraining still happens at 8K-scale lengths before a dedicated long-context stage. Deploying a model on inputs longer than it was trained for requires \textit{context extension}: either modifying the model at inference time so it generalizes beyond its training length (this section), or training it for the target length directly (Section~\ref{subsec:native_long_context}, which is how frontier models actually get their 128K--1M windows).

The bottleneck is almost always \textbf{positional encoding}. Absolute learned embeddings cannot extrapolate at all. Sinusoidal encodings degrade quickly. The modern approach is Rotary Position Embeddings (RoPE), which encodes relative position by rotating query and key vectors in 2D subspaces.

\subsection{RoPE and Why It Enables Extension}

\textbf{What it is.} RoPE applies a rotation matrix $\mR_{\theta, m}$ to query and key vectors at position $m$. The attention dot product between positions $m$ and $n$ depends only on the relative distance $m - n$, not on absolute position.

\begin{mathresult}{RoPE Attention}
\[
\text{Attention}(m, n) = (\mR_{\theta, m}\,\vq)^\top (\mR_{\theta, n}\,\vk) = \vq^\top \mR_{\theta, n-m}\,\vk
\]
where $\mR_{\theta, m}$ rotates each pair of dimensions $(2i, 2i{+}1)$ by angle $m \cdot \theta_i$, with base frequencies $\theta_i = \text{base}^{-2i/d}$ and $\text{base} = 10{,}000$ by default.
\end{mathresult}

\textbf{Why extension is possible.} Because RoPE encodes position through rotation angles, you can manipulate the frequencies to ``compress'' longer sequences into the model's trained range. The key parameter is the \textbf{base frequency}---changing it changes the wavelength of each rotary dimension.

\subsection{Inference-Time Extension Methods}
\label{subsec:inference_time_extension}

The methods in this table are what you reach for when you \textit{cannot retrain}: they rescale RoPE at inference time (optionally with a brief fine-tune) so that positions beyond the training length map back into the trained distribution. They dominated 2023 practice and remain the standard emergency path for a short-trained model---but they are retrofits. Frontier models do not reach 128K--1M this way; they are trained for length (Section~\ref{subsec:native_long_context}). Position the two correctly in an interview: this table is the ``ship next month'' answer, the training recipe is the ``how it is actually done'' answer.

\begin{table}[H]
\centering
\caption{RoPE-based context extension methods (inference-time retrofits, optionally with brief fine-tuning)}
\begin{tabularx}{\textwidth}{lXXc}
\toprule
\textbf{Method} & \textbf{Key Idea} & \textbf{Trade-off} & \textbf{Year} \\
\midrule
Position Interpolation & Linearly scale positions by $\frac{L_{\text{train}}}{L_{\text{target}}}$; all frequencies compressed equally & Degrades high-frequency (local) positional info; needs fine-tuning & 2023 \\
NTK-aware Scaling & Increase RoPE base from 10K to a higher value; spreads compression across frequencies nonuniformly & Better preserves local position; minimal fine-tuning needed & 2023 \\
YaRN & Combine NTK scaling with temperature correction and attention scaling factor & Best quality at high extension ratios (16--32$\times$); needs brief fine-tuning & 2023 \\
Dynamic NTK & Adjust the base frequency dynamically at inference based on current sequence length & No fine-tuning needed; quality degrades at extreme lengths & 2023 \\
\bottomrule
\end{tabularx}
\end{table}

\textbf{Position Interpolation} divides all position indices by a scaling factor $s = L_{\text{target}} / L_{\text{train}}$. A model trained on 4K, extended to 16K, maps position 16,000 to effective position 4,000. Simple but treats all RoPE frequency bands identically, which is suboptimal---low-frequency bands (capturing global position) need more compression than high-frequency bands (capturing local position).

\textbf{NTK-aware scaling} modifies the RoPE base frequency: instead of $\text{base} = 10{,}000$, use $\text{base}' = 10{,}000 \cdot s^{d/(d-2)}$. This applies stronger compression to low-frequency bands and weaker compression to high-frequency bands, preserving local positional resolution. The name comes from the analogy to Neural Tangent Kernel theory, where high-frequency components require finer resolution.

\textbf{YaRN} (Yet another RoPE extensioN) is the most complete approach. It combines NTK-aware base scaling with a per-dimension interpolation strategy (dimensions below a threshold are not scaled at all), plus a learned or computed attention scaling factor ($\sqrt{t}$) that corrects for the entropy increase caused by longer sequences.

\begin{keyinsight}[Why Naive Extrapolation Fails]
A model trained on 4K tokens has never seen RoPE angles corresponding to position 5,000+. These out-of-distribution angles produce attention patterns the model was never trained to interpret---attention scores become essentially random beyond the training length. Context extension methods work by ensuring all position encodings at the target length map back into the distribution of encodings the model saw during training. The choice of method determines how gracefully this mapping preserves local vs.\ global positional information.
\end{keyinsight}

\subsection{Non-RoPE Approaches}

Not all long-context solutions modify position encodings:

\begin{itemize}
    \item \textbf{ALiBi} (Attention with Linear Biases): Adds a position-dependent penalty directly to attention logits. No learned positional parameters. Extrapolates naturally but cannot match RoPE quality at very long contexts.
    \item \textbf{Ring Attention}: Distributes the long sequence across multiple devices, with each device computing attention for its local chunk and passing KV blocks in a ring topology. Enables context lengths limited only by the number of devices, not memory.
    \item \textbf{Landmark / infini-attention}: Compress old context into fixed-size memory states, maintaining a ``compressive memory'' that the model can attend to. Trades exact retrieval for bounded memory.
    \item \textbf{Training on long data}: The most reliable method, and how every current frontier model (128K--1M+ windows) actually gets its context length. Avoids all extension artifacts at the cost of a dedicated training stage---covered next.
\end{itemize}

\subsection{How Long Context Is Actually Achieved (2024--2026)}
\label{subsec:native_long_context}

The PI/NTK/YaRN table describes retrofits. Since 2024, frontier practice treats long context as a \textbf{training recipe}: a continued-pretraining stage with progressive length extension and a long-data mixture, run between the main pretraining run and post-training. The Llama~3.1 128K recipe is the canonical published shape:

\begin{enumerate}
    \item \textbf{Pretrain at standard length.} The bulk of tokens (trillions) are trained at 8K. Long sequences are expensive---attention FLOPs grow quadratically and activation memory linearly---and most of what the model learns does not require them.
    \item \textbf{Raise the RoPE base.} Before the long stage, increase the base frequency (Llama~3.1 uses 500{,}000 instead of 10{,}000) so the lowest-frequency dimensions have wavelengths that span the target window. This is the same lever as NTK-aware scaling, but applied \textit{before training on long data} rather than as an inference-time patch.
    \item \textbf{Continued pretraining with progressive length extension.} Increase the context length in stages (Llama~3.1: six stages from 8K to 128K, $\sim$800B tokens total---roughly 5\% of the full token budget). At each stage, train until long-context probes (needle-in-a-haystack) recover \textit{and} short-context benchmarks have not regressed, then lengthen again.
    \item \textbf{Long-data mixture.} Upsample naturally long documents (books, code repositories, long web pages) but keep the mixture mostly short. Training exclusively on long data degrades short-context quality; the mixture ratio is a tuned hyperparameter, not an afterthought.
    \item \textbf{Long-context post-training.} Include long-context SFT data (long-document QA, summarization, repo-level code tasks---largely synthetic) so instruction following survives at length. Skipping this step produces a model that passes base-model needle tests but ignores instructions at 100K.
\end{enumerate}

Two named 2024 methods extend this picture, one on each side of the train/serve divide:
\begin{itemize}
    \item \textbf{LongRoPE}: replaces hand-derived rescaling schedules (NTK, YaRN) with an evolutionary \textit{search} over non-uniform per-dimension interpolation factors, combined with progressive extension---reaching 2M-token windows from short-context models with modest fine-tuning.
    \item \textbf{Dual-chunk attention (DCA)}: a \textit{training-free} serving-time method that splits positions into chunks and remaps intra-chunk vs.\ inter-chunk relative positions so no rotary angle ever exceeds the trained range (used in Qwen's long-context serving).
\end{itemize}

\textbf{Infrastructure corollary.} Training at 128K+ does not fit a single GPU's activation memory: the long stage requires \textbf{context parallelism}---sharding the sequence dimension across devices with ring-style KV exchange, the training-side counterpart of the Ring Attention idea above. See Chapter~\ref{chap:training_optimization} for the parallelism taxonomy.

\begin{keyinsight}[Extension Tricks vs.\ the Native Recipe]
When an interviewer asks ``how would you give this model 128K context?'', the 2026 answer starts with the continued-pretraining recipe (raise the base, progressive lengths, long-data mixture, long SFT) and reaches for PI/NTK/YaRN only under the constraint ``you cannot retrain and must ship now.'' Describing YaRN as the state of the art---rather than as the best of the retrofits---dates you to 2023.
\end{keyinsight}

\subsection{Evaluating Long Context}
\label{subsec:long_context_eval}

Every model advertises a \textbf{claimed} context length---the maximum positions it accepts. What matters is the \textbf{effective} context length: the longest input at which it still does the task. The two routinely diverge, and ``the model claims 1M tokens---how do you verify its effective context?'' is now a standard interview question.

\textbf{Needle-in-a-haystack (NIAH), and why it saturated.} Insert a short out-of-place fact (the ``needle'') at controlled depths in long distractor text and ask the model to retrieve it; plot accuracy as a depth $\times$ length heatmap. NIAH was the de facto long-context eval of 2023--2024 and is now a \textit{saturated smoke test}: retrieving one verbatim fact is far easier than using a long context, and frontier models pass NIAH at 1M tokens while failing aggregation or multi-hop tasks at 64K. Passing NIAH is necessary, not sufficient---a red-green NIAH grid still catches position-OOD cliffs and serving bugs, which is why it survives as a regression test.

\textbf{The harder eval classes:}
\begin{itemize}
    \item \textbf{RULER-style synthetic suites}: scale difficulty past single-needle retrieval---multiple needles, multi-hop variable tracking (follow a chain of aliased variables across the context), and \textit{aggregation} (e.g., extract the most frequent words), plus long QA. RULER defines effective context length as the longest input at which the score stays above a fixed baseline; its headline 2024 finding was that roughly half of the models claiming 32K+ failed to hold threshold performance even at 32K, and claimed-vs-effective gaps of $4\times$ or more remain routine.
    \item \textbf{Multi-hop reasoning}: the answer requires chaining facts located at distant positions; degrades much earlier than single-fact retrieval.
    \item \textbf{Aggregation / holistic tasks}: summarization, counting, theme extraction over the \textit{entire} input. The hardest class and where degradation appears first, because no retrieval shortcut exists.
    \item \textbf{Long-CoT / long generation}: reasoning models must also \textit{produce} long outputs; evaluating reasoning quality over a long deliberation is a separate axis from reading length, and realistic suites (LongBench-style long-document QA, repo-level code tasks) increasingly combine both.
\end{itemize}

\textbf{Lost in the middle, quantified.} On multi-document QA, Liu et al.\ measured accuracy dropping by up to $\sim$20 points when the answer-bearing document sits mid-context rather than first or last---a U-shaped curve over answer position. The effect is \textit{partially fixed by training}: models that went through a native long-context stage flatten the U considerably at moderate lengths, but the dip re-emerges as you approach a model's effective-length limit. Treat position sensitivity as a property to measure, not an assumption in either direction.

\begin{warningbox}
\textbf{Never accept a claimed context length.} Before serving a long-context workload, run your own grid: NIAH as the smoke test, then a RULER-style aggregation/multi-hop suite at \textit{your} target lengths, with answer positions permuted. Extension retrofits in particular can hold perplexity flat while retrieval-in-context collapses---perplexity is not a long-context eval.
\end{warningbox}

%==============================================================================
\section{The RAG Pipeline in Brief}
\label{sec:rag_architecture}
%==============================================================================

\textit{Pointer: the RAG pipeline is canonical in the NLP volume [NLP~9]---the offline/online architecture in full, chunking strategy, retrieval-paradigm selection, query transformation, grounded generation and citation, the advanced patterns (Self-RAG, agentic RAG, context compression), RAG evaluation, and the failure taxonomy that says which layer broke. The retrieval infrastructure beneath it is Volume~III: vector-index internals and sizing [SR~3], bi-encoder and reranker training and hybrid fusion [SR~4], and index freshness, ANN serving, caching, and monitoring [SR~8]. This section keeps only what the context-versus-retrieval decision of Section~\ref{sec:long_context_vs_rag} requires: the shape of the system you would be committing to build, what it costs to own, and how it fails.}

\textbf{The pipeline in one paragraph.} Offline, documents are parsed, split into chunks, embedded, and written to a vector index alongside their text and metadata, with a lexical index built over the same chunks. Online, the query is optionally rewritten or decomposed, run against both indexes, the two candidate lists are fused, the top 20--100 candidates are reranked, and the surviving 3--10 chunks are assembled into a grounded prompt (``answer only from the provided context; cite sources; say you do not know if the context is insufficient''). Two-stage retrieval exists because the cheap scorer is decomposable---document representations are precomputed, so first-stage search is sublinear in corpus size---while the accurate scorer needs the query and document together and therefore runs only over a short candidate list. The invariant that matters for the architectural decision: \emph{a chunk retrieval misses is a fact the generator does not have}, and nothing downstream recovers it.
\label{fig:rag_pipeline}% label retained; the full pipeline diagram is canon in [NLP 9]

\subsection{Chunking}
\label{sec:chunking}

Chunking fixes the granularity at which anything can ever be retrieved. The menu---fixed-size with overlap, recursive/structural, semantic, document-aware, and the parent-child pattern that indexes small chunks for precision but returns their parent section for generation---is canonical in [NLP~9], along with how to choose among them. Three facts carry into this chapter's decision. (1)~\textbf{Defaults}: 256--512 tokens with 10--20\% overlap, never splitting mid-sentence; treat chunk size as a hyperparameter tuned on a retrieval eval set, not a constant. (2)~\textbf{The chunk multiplier is the real corpus size}: 10M documents at roughly five chunks each is a 50M-vector index, and that multiplier is re-paid on every re-embedding pass ([SR~3] for the sizing arithmetic). (3)~\textbf{Chunking is a lossy, irreversible index-time commitment}: it bounds what the retriever can ever see, and revising it means re-indexing the corpus---which is precisely the class of decision a context window lets you skip entirely.

\subsection{Query Processing}
\label{sec:query_processing}

Raw user queries are poor retrieval probes: short, conversational, and written in the user's vocabulary rather than the corpus's. The three standard transforms---\textbf{rewriting} (resolve pronouns and acronyms against conversation history), \textbf{HyDE} (embed an LLM-generated hypothetical answer instead of the question, bridging query--document style mismatch), and \textbf{decomposition} (split a multi-part question into sub-queries retrieved independently, then synthesize)---are canonical in [NLP~9], with their failure modes and a routing framework for choosing among them. What matters here is the cost line they add: every transform spends LLM calls \emph{before} retrieval even begins, so a ``cheap'' RAG query is rarely a single model call, and a decomposed multi-hop query is $N$ retrievals plus $N{+}1$ generations---the regime where one long-context pass stops looking expensive (Section~\ref{sec:long_context_vs_rag}).

\subsection{Reranking and Hybrid Retrieval}
\label{sec:reranking}

\textit{Pointer: the retrieval-systems depth behind this section---bi-encoder training recipes, ColBERTv2/PLAID and the late-interaction storage math, LLM rerankers and their distillation, fusion mechanics, and multi-stage latency/cost budgets---is canonical in the Search volume [SR~4]; the embedder landscape is canonical in [NLP~3]. This section keeps what a RAG designer needs to assemble the pipeline.}

\textbf{The funnel in one paragraph.} Retrieval and reranking form a funnel that spends more compute per candidate on fewer candidates. The first stage casts a wide net with \textit{decomposable} scorers whose document representations are precomputed offline: BM25 and a bi-encoder run in parallel, and their candidate lists are merged with reciprocal rank fusion (rank-based fusion---raw BM25 and cosine scores live on incompatible, query-dependent scales and must not be summed). The hybrid matters because the two retrievers fail on complementary queries: embedding models blur exact identifiers (product IDs, error codes, proper nouns) that BM25 matches literally, while BM25 misses the paraphrases that embeddings catch---so hybrid consistently outperforms either alone, and in an interview you should always propose it. A \textbf{cross-encoder} then rescores the top 20--100 candidates jointly (\texttt{[CLS] query [SEP] chunk [SEP]}): full query--document attention buys back the precision that single-vector similarity loses, at $O(N)$ forward passes with nothing precomputable---affordable only over a small window, never over the corpus. \textbf{Late interaction} (ColBERT) sits between the two: per-token document embeddings stored at index time, scored by MaxSim at query time---near-cross-encoder quality with bi-encoder-style precomputation, at a storage multiple (roughly an order of magnitude over a single-vector index even with ColBERTv2's compression; the worked math is in [SR~4]).

\textbf{Practical default for RAG:} hybrid BM25 + bi-encoder retrieval fused with RRF, cross-encoder reranking of the top 20--50, and the top 3--10 chunks to the generator---where the reranker's output \textit{ordering}, not just its selection, affects generation quality (Section~\ref{sec:long_context_vs_rag}). When quality demands more, the modern upgrade path is distilling an LLM listwise reranker into the serving cross-encoder rather than serving anything larger online [SR~4].

\subsection{What the Pipeline Costs to Own}

Per-query token cost is the visible half of the comparison with long context (the arithmetic is in Section~\ref{sec:long_context_vs_rag}). The other half is standing cost that a long prompt does not incur at all, and it is what teams systematically under-count:

\begin{itemize}
    \item \textbf{Components to run.} An ingestion and parsing pipeline per document format, an embedding service, a vector index, a lexical index, a fusion and rerank tier, and prompt assembly. Each is a deployment, a dashboard, and an on-call surface ([SR~8]).
    \item \textbf{Standing infrastructure.} Index memory is billed hourly whether or not queries arrive, and in a served RAG stack the reranking GPUs---not the vector search---are typically the dominant line item; the worked fleet economics are [SR~8].
    \item \textbf{The cost of change.} Re-embedding is the unit of change: a new embedding model, chunking rule, or parser means a full corpus pass and a dual-index migration. Embedding model and index artifact are versioned together or retrieval silently mixes vector spaces.
    \item \textbf{Freshness.} RAG's product claim is knowledge past the training cutoff, so index lag re-imposes a cutoff and the failure mode is a fluent, well-cited, stale answer ([SR~8]).
    \item \textbf{Evaluation.} A labeled (query, relevant-chunk) set and a faithfulness harness must exist and be maintained, or none of the above can be tuned---and the labeling burden is permanent, because the corpus changes.
    \item \textbf{Access control.} Retrieval is where per-user permissions have to be enforced, since anything retrieved can leak into the answer; ACLs are pre-filters, not post-filters ([SR~8]).
\end{itemize}

None of this argues against RAG: above a certain corpus size it is the only option. It argues that the comparison in Section~\ref{sec:long_context_vs_rag} is between a \emph{system} and a \emph{prompt}, and that counting tokens alone flatters the system.

%==============================================================================
\section{Beyond Top-$k$: GraphRAG, Contextual Retrieval, Agentic Retrieval}
\label{sec:modern_rag}
%==============================================================================

The pipeline above---chunk, embed, retrieve top-$k$, rerank, generate---is the 2023 canonical stack. Three 2024--2026 developments extend it in different directions: GraphRAG changes \textit{what is indexed}, contextual retrieval and late chunking change \textit{how chunks are represented}, and agentic retrieval changes \textit{who drives the loop}.

\subsection{GraphRAG}

\textit{Pointer: the core RAG pipeline that GraphRAG extends is treated canonically in the NLP volume [NLP~9]; this subsection covers only the graph-based extension.}

\textbf{The structural failure of top-$k$.} Vector retrieval answers \textit{local} questions---ones whose answer lives in a few passages. It structurally cannot answer \textit{corpus-level} questions: ``What are the main themes across these 10{,}000 support tickets?'' or ``How did our contract terms evolve over five years?'' No top-$k$ set of chunks contains the answer, because the answer is a property of the corpus, not of any passage. Retrieving more chunks does not fix this; the information must be \textit{aggregated at index time}.

\textbf{Construction (Microsoft GraphRAG, 2024).} An LLM pass over every chunk extracts entities and relations, which are assembled into a knowledge graph. Community detection (Leiden) then partitions the graph into a hierarchy of communities, and the LLM writes a \textbf{community summary} for each node at each level---pre-computed aggregations of the corpus at multiple granularities.

\textbf{Query time.} Route by query type:
\begin{itemize}
    \item \textbf{Global queries} (themes, trends, corpus-wide comparisons): map-reduce over community summaries---each relevant summary produces a partial answer, a final LLM call synthesizes them. This is what vanilla RAG cannot do at all.
    \item \textbf{Local queries} (about a specific entity): enter the graph at the matched entities, expand along edges to neighbors and their source chunks. This also improves multi-hop entity questions, since relations are explicit edges rather than embedding-space neighbors.
\end{itemize}

\textbf{Cost and when to use it.} Index construction runs LLM calls over the entire corpus---orders of magnitude more expensive than embedding, and incremental updates are awkward (entity extraction is cheap to append; community structure and summaries need periodic rebuilds). Use GraphRAG when corpus-level sensemaking or entity-centric multi-hop questions are a first-class requirement; stay with vector top-$k$ for straightforward factual QA.

\subsection{Contextual Retrieval and Late Chunking}

\textit{Pointer: chunking fundamentals (strategies, sizes, parent-child) are covered above and canonically in [NLP~9]; this subsection covers the 2024 fix for context-starved chunks.}

\textbf{The problem: chunks lose their document.} A chunk reading ``The company's revenue grew 3\% over the previous quarter'' is nearly useless in isolation---which company? which quarter? Pronouns, dangling references, and missing section headers make chunk embeddings poorly discriminative and defeat BM25 entirely. Parent-child retrieval fixes the \textit{generation} side (return more context to the LLM) but not the \textit{retrieval} side: the small chunk is still embedded without its context.

\textbf{Contextual retrieval (Anthropic, 2024).} At index time, prompt an LLM with the \textit{full document plus the chunk} to generate 50--100 tokens of situating context (``This chunk is from ACME Corp's Q2 2023 SEC filing; the previous quarter's revenue was \ldots''), and prepend it to the chunk before embedding \textit{and} before BM25 indexing (contextual embeddings + contextual BM25). Anthropic reports top-20 retrieval failure rate reduced by $\sim$35\% with contextual embeddings, $\sim$49\% combined with contextual BM25, and $\sim$67\% with reranking added. The cost---one LLM call per chunk at index time---is made tolerable by prompt caching: the full-document prefix is cached once and reused across all of that document's chunks, so each chunk pays only for its own tokens plus the cached-input rate.

\textbf{Late chunking.} A generation-free alternative: run the \textit{full document} through a long-context embedding model in one pass, then mean-pool the contextualized token embeddings within each chunk boundary \textit{afterwards}. Every chunk embedding is conditioned on the whole document (the tokens attended across chunk boundaries before pooling) without any LLM calls. Requires a long-context embedding model and only helps the dense side.

\textbf{Choosing between them.} Contextual retrieval costs LLM tokens at index time but works with any embedder and fixes BM25 too; late chunking costs one embedding pass and nothing else but leaves sparse retrieval context-blind. Both compose with parent-child retrieval, which remains the generation-side complement.

\subsection{Agentic Retrieval}

\textit{Pointer: agentic RAG as a \emph{pattern}---retrieval as a tool the model may invoke zero or more times, multi-hop loops, and Self-RAG/CRAG-style reflection that decides whether to retrieve at all and grades what came back---is canonical in [NLP~9]; the serving consequences of an agent issuing tens of retrieval calls per task are [SR~8], and agent architecture is Chapter~\ref{chap:agents_tool_use}. What belongs here is the architectural question the loop raises: when it replaces the index.}

\textbf{From pipeline to loop.} The structural fact to carry forward is that a loop conditions each query on the previous hop's results, so multi-hop questions are answered natively---at the price of several sequential LLM calls, and therefore seconds, per user question.

\textbf{The live debate: does an agent with grep beat a vector DB?} Coding agents demonstrated that an agent iterating with exact-match search (\texttt{grep}/ripgrep) and file reads over a raw corpus often outperforms one-shot vector retrieval---on codebases, decisively. The reasons generalize: queries against code and structured corpora are precise identifiers (exact-match wins, embeddings paraphrase); the agent \textit{sees its results and iterates}, so a bad first query costs one loop iteration rather than the answer; and no index means no staleness, no chunking decisions, and no embedding-drift maintenance. The vector DB retains the advantage when the corpus is huge (iterating with grep over 25B tokens does not converge), when vocabulary mismatch is the dominant failure mode (users paraphrase; documents do not), and when a sub-second, fixed-latency single pass is required---an agentic loop costs several LLM calls and seconds to minutes. The 2026 trend is that cheaper, faster models plus prompt caching keep moving the crossover toward agentic search; the defensible interview answer is a hybrid: semantic retrieval as \textit{one tool among several} (alongside grep, metadata filters, and a citation graph) inside an agent loop, with the loop budget matched to query value.

\begin{keyinsight}[Retrieval Is Becoming a Tool, Not a Pipeline]
The 2023 question was ``design a RAG pipeline.'' The 2026 question is ``which retrieval tools do you expose to the agent, and when does the agent skip the vector DB entirely?'' Answering the second question well requires the first---the pipeline components do not disappear; they become tools whose invocation the model controls.
\end{keyinsight}

%==============================================================================
\section{Where RAG Breaks}
\label{sec:rag_eval}
%==============================================================================

\textit{Pointer: RAG evaluation is canonical in [NLP~9]---the three-layer decomposition (retrieval quality, generation quality, faithfulness), the RAGAS metric set and how an LLM judge computes each, evaluation-set construction, and the full failure taxonomy. Retrieval metric definitions are [SR~7]; monitoring a live retrieval stage is [SR~8]. What this chapter keeps is the failure \emph{surface}, because half of the context-versus-retrieval decision is which failures you can afford.}

A RAG answer fails in four distinguishable ways, and the differences from long-context failure are what make the decision:

\begin{enumerate}
    \item \textbf{Retrieval miss}---the needed chunk is not in the top-$k$. Unrecoverable downstream and the most common root cause by a wide margin; check recall@$k$ on a sample of failures before touching prompts or models. Measured by recall@$k$ and context recall.
    \item \textbf{Ranking failure}---the chunk is retrieved but ranked low, landing where the generator under-attends. Ordering, not just selection, is what reaches the model.
    \item \textbf{Grounding failure}---the right context is present and the answer is still wrong: the model answers from parametric memory, drowns in too many chunks, or hits a conflict between context and priors. Measured by faithfulness against the retrieved context.
    \item \textbf{Chunk-boundary failure}---the answer spans material that chunking separated and retrieval never returned together. This is the failure mode that most cleanly argues for a longer window: the fix is more context, not better ranking.
\end{enumerate}

Component metrics do not compose. A better retriever can surface chunks that confuse the generator, so component gains must be confirmed end-to-end; and faithfulness can be high precisely because the model correctly refused to answer from a bad context, so faithfulness must always be read next to context recall.

\begin{keyinsight}[The Two Architectures Fail on Different Axes]
A long-context system fails by \emph{position and length}: degradation toward the middle of the window and a cliff past the effective length, both visible on a NIAH/RULER length$\times$depth grid (Section~\ref{subsec:long_context_eval}). A RAG system fails by \emph{recall}: whatever the retriever did not return is invisible to every downstream measurement, so no amount of generator evaluation can see it. The instruments are therefore different---length--depth grids versus recall@$k$ and faithfulness---and a hybrid needs both, because retrieving 30 chunks into a 128K window buys cross-chunk reasoning while still inheriting the retriever's recall ceiling.
\end{keyinsight}

%==============================================================================
\section{Long Context vs RAG}
\label{sec:long_context_vs_rag}
%==============================================================================

This is the decision the chapter exists for, and the one the rest of the program points here to make: [NLP~9] frames the same trade from the pipeline builder's side (given that you are shipping RAG, where is retrieval the wrong tool?) and [SR~8] from the serving side. The naive version---``windows are 1M tokens now, can we delete the retriever?''---has a short answer and a long one. Short: no, because corpora exceed any window and cost scales with what you put in it. Long: the boundary moves every year, and a Staff-level answer says what moves it.

\textbf{Step 1: size the corpus against the window.} Only one of the three regimes is a real decision.
\begin{itemize}
    \item \textbf{Corpus $\ll$ window} (one contract, one repository, a day of logs): stuff it. Building retrieval here is unpaid complexity; the only open questions are cost per call and whether prefix caching applies.
    \item \textbf{Corpus $\gg$ window} (10K+ documents, any enterprise wiki---25B tokens is a routine number): retrieve. No window is coming that holds it, and you would not pay to prefill it per query if one were.
    \item \textbf{Corpus within an order of magnitude of the window}: the genuine trade, decided by task shape, per-query economics, and freshness.
\end{itemize}

\textbf{Step 2: check the window you actually have, not the one on the model card.} A 128K claim buys nothing if effective length ends at 32K, and the gap is routinely $4\times$ or more (Section~\ref{subsec:long_context_eval}). Run the length$\times$depth grid at your target length \emph{before} the architecture argument, because it can settle it.

\begin{table}[H]
\centering
\caption{Long context vs.\ RAG: when to choose each}
\begin{tabularx}{\textwidth}{lXX}
\toprule
\textbf{Dimension} & \textbf{Long Context (stuff everything)} & \textbf{RAG (retrieve then generate)} \\
\midrule
Knowledge base size & Up to context limit (128K--1M tokens) & Unlimited (millions of documents) \\
Latency & Higher (process entire context) & Lower (process only retrieved chunks) \\
Cost per query & High (billed per input token) & Lower (fewer input tokens) \\
Accuracy (needle-in-haystack) & Degrades in middle of long contexts & Depends on retrieval quality \\
Freshness & Must re-inject updated documents & Update index; no model change \\
Setup complexity & Low (just concatenate) & High (chunking, embedding, indexing, reranking) \\
Reasoning across documents & Strong (all docs visible simultaneously) & Weak (only sees retrieved subset) \\
Failure mode & ``Lost in the middle'' attention degradation & Retrieval misses relevant documents \\
Attribution & Possible but unenforced (cite by passage) & Native: chunk $\to$ document and page \\
Access control & All-or-nothing---whatever is in the prompt is visible & Enforced at retrieval, as a pre-filter [SR~8] \\
Standing cost & None beyond the model & Index memory, rerank fleet, freshness and evaluation pipelines \\
\bottomrule
\end{tabularx}
\end{table}

\textbf{Step 3: choose on task shape.} Long context wins when the task is \emph{holistic}---summarize the whole deal, find contradictions across a contract, count or aggregate over the corpus, reason across documents that no query would have co-retrieved---because such answers are not localized in any top-$k$ set and retrieving more chunks does not converge on them (the index-time alternative is aggregation, Section~\ref{sec:modern_rag}). RAG wins when answers are \emph{localized} and the corpus is large, fresh, or permission-scoped: attribution is native, updates are an index write rather than a re-upload, and cost per query is one to two orders of magnitude lower.

\textbf{Step 4: price it, with caching in the model.} At illustrative rates (\$10/M input tokens, cached input at $\sim$0.1$\times$; state prices as variables, they fall every year), a full 128K window is $\sim$\$1.28 per query against $\sim$\$0.04 for RAG over five chunks---roughly $30\times$. Cached input changes the comparison only when the \emph{same prefix repeats}: a multi-turn session over one document set amortizes to $\sim$5--6$\times$, close enough that quality should decide; independent queries over a shifting corpus never hit the cache and stay at the full multiple. The worked arithmetic, including prefill and time-to-first-token, is the cost question in Section~\ref{sec:rag_interview}.

\textbf{Step 5: remember the serving bill is not linear in the window.} Prefill compute grows with the quadratic attention term and KV-cache bytes grow linearly with length, so a long window degrades concurrency and time-to-first-token, not just the invoice (Chapter~\ref{chap:inference_optimization}). RAG's per-query serving profile is flat by construction: the generator always sees a few thousand tokens.

\textbf{What moves the crossover.} Four forces, and they do not all push the same way. (1)~\textbf{Cheaper long context}---falling prices and cached input push toward the window for repeated-prefix workloads. (2)~\textbf{Better long-context training}---native long-context stages flatten the lost-in-the-middle curve and raise effective length, making stuffing more reliable. (3)~\textbf{Better retrieval}---hybrid retrieval, rerankers, and contextual retrieval raise recall, which is RAG's ceiling; when recall@$k$ is already 0.95 a bigger window buys little, and when it is 0.6 the window is a genuine quality upgrade. (4)~\textbf{Agentic search}---an iterating agent turns a recall problem into a latency-and-cost problem, and on precise, identifier-heavy corpora it can beat one-shot vector retrieval outright (Section~\ref{sec:modern_rag}). The 2026 shape of the answer is that the extremes are settled and the middle is a routing decision, not an architecture religion.

\textbf{The routing rule, workload by workload.} In production the answer is rarely global; it is a router.

\begin{table}[H]
\centering
\caption{Which architecture each workload wants}
\begin{tabularx}{\textwidth}{lXX}
\toprule
\textbf{Workload} & \textbf{Choice} & \textbf{Why} \\
\midrule
Analyze one long document (contract, filing, incident log) & Long context & Corpus fits; the task is holistic; chunking would destroy the cross-references \\
Corpus-wide sensemaking (themes across 10K tickets) & Neither alone & No top-$k$ set contains the answer and no window holds the corpus---aggregate at index time (Section~\ref{sec:modern_rag}) \\
Factual lookup over a large, fresh corpus & RAG & Localized answer, attribution required, index updates in minutes \\
Multi-hop entity question & RAG in a loop & The second query must be conditioned on the first hop's result \\
Code navigation over a repository & Agentic exact-match search & Identifiers favor literal matching; the agent iterates and sees its own results \\
Multi-turn session over a fixed document set & Long context with prefix caching & Cached input collapses the repeated prefix to $\sim$0.1$\times$ and skips prefill \\
Permission-scoped enterprise Q\&A & RAG & ACLs are enforced at retrieval; a prompt has no access control \\
\bottomrule
\end{tabularx}
\end{table}

\textbf{The hybrid's operating parameters.} ``Retrieve wide into a long window'' is a real configuration, not a slogan: retrieve 20--30 chunks rather than 3--5 because the window can hold them, rerank to order rather than to prune, place the highest-scoring chunks first and last, and keep the assembled context in the 10K--30K range where measured degradation is mild for a natively long-context model. Escalate to the full window only when a holistic pass is genuinely required, and de-escalate---fewer, better chunks---when faithfulness rather than recall is the failing metric. The two knobs to tune with evidence are the retrieval depth (chosen from the recall curve at your target latency) and the assembled-context length (chosen from your own length$\times$depth grid, not from the model card).

\begin{warningbox}
\textbf{``Lost in the middle'' is a real production risk.} Liu et al.\ (2024) showed that LLMs pay most attention to the beginning and end of their context window: on multi-document QA, accuracy drops by up to $\sim$20 points when the answer-bearing document sits mid-context rather than first or last---a U-shaped curve that native long-context training flattens but does not eliminate (Section~\ref{subsec:long_context_eval}). For long-context approaches, order your documents with the most relevant first and last. For RAG, this means the reranker's ordering directly impacts generation quality---do not treat the top-$k$ results as an unordered set.
\end{warningbox}

\begin{keyinsight}[Choose by Which Failure You Can Afford]
Cost picks the default and task shape overrides it, but the durable framing is failure mode. A long-context system loses information \emph{in the middle of its window and past its effective length}; a RAG system loses \emph{whatever the retriever did not return}, invisibly. Holistic tasks have no retrievable answer and must take the window (or index-time aggregation); localized questions over a large, fresh, access-controlled corpus have no affordable window and must take retrieval. The production default in between is neither: retrieve wide (20--30 chunks, 10K--30K tokens), rerank hard, order the survivors with the best first and last, and let a long-context model reason over all of them---you pay one prefill you can afford and inherit the recall ceiling without the lost-in-the-middle curve. If you can only say one sentence: \emph{context is where you put what you already know is relevant; retrieval is how you find out.}
\end{keyinsight}

%==============================================================================
\section{Interview Questions}
\label{sec:rag_interview}
%==============================================================================

\begin{keyinsight}[Where the RAG Questions Live]
This bank is deliberately narrow: long-context extension, long-context debugging, and the context-versus-retrieval trade. The RAG \emph{pipeline} questions that used to sit here are canonical in [NLP~9] and answered there in full---end-to-end RAG system design (and its 10M-document, 1000-QPS scaling variant), chunking strategy and chunk-size choice, HyDE versus other query transformations, RAG evaluation end-to-end, ``retrieval is right but the answer is wrong'' debugging (including the multi-hop and reasoning-query variants, whose fixes---decomposition, multi-hop retrieval, agentic loops---are that chapter's advanced-patterns material), lost-in-the-middle handling, bi-encoder versus cross-encoder with ColBERT as the middle ground, and RAG versus fine-tuning versus long context from the pipeline builder's side. Multi-stage retrieval design as a \emph{retrieval-system} question, with its recall and latency budgets, is [SR~4]; scaling and operating a live retrieval stage is [SR~8]. Prepare those there. Prepare the four below here---they are the ones an interviewer asks when the subject is the model's window rather than the pipeline around it.
\end{keyinsight}

%--- Q1: Long-context extension under deadline ---

\begin{interviewq}
\textbf{Q: Your 8K-trained model must serve 64K contexts next month. What are your options, and what do they cost?}

\badgeLsix{L6} \quad \badgetype{Trade-off} \quad \badgefreq{\freqhigh}

\stronganswer

\textbf{Frame the option space first.} An $8\times$ extension with a one-month deadline puts five options on the table, at different points on the quality--effort--risk frontier:

\textbf{Option 1: Training-free RoPE rescaling (Dynamic NTK, dual-chunk attention).} Shippable this week: change the RoPE base (or chunk the position mapping) at inference so no rotary angle leaves the trained range. Zero training cost. Quality degrades toward the far end of the window, and $8\times$ is near the edge of what training-free methods deliver---acceptable as a stopgap behind a flag, not as the endpoint.

\textbf{Option 2: PI / NTK-aware / YaRN with a brief fine-tune.} YaRN at $8$--$16\times$ with a short fine-tune on long sequences (a few hundred million to a few billion tokens, a tiny fraction of pretraining) is the best quality-per-GPU-hour among the retrofits (Section~\ref{subsec:inference_time_extension}). Fits the month, including data prep and evaluation. Risks: fine-tuning only on long data regresses short-context quality (mix lengths), and perplexity can stay flat while retrieval-in-context fails---gate the launch on NIAH plus a RULER-style suite at 64K, not on loss curves (Section~\ref{subsec:long_context_eval}).

\textbf{Option 3: Continued pretraining with progressive lengths (the frontier recipe).} Raise the RoPE base, extend in stages 8K$\to$16K$\to$32K$\to$64K on a long-data mixture, then long-context SFT (Section~\ref{subsec:native_long_context}). Most robust; also the most expensive: long-document data curation, context-parallel training infrastructure (Chapter~\ref{chap:training_optimization}), and multiple weeks. The right roadmap item; risky as the only plan for a one-month deadline.

\textbf{Option 4: Don't extend---RAG over the long input.} If ``64K support'' really means ``answer questions over long documents,'' chunk and retrieve within the existing 8K window. No model work at all. Fails when the task is holistic (summarize the whole deal, find contradictions across the document), which is exactly where long context earns its cost.

\textbf{Option 5: Buy it.} Serve a model already trained to 64K+ (API or open weights) for the long-context traffic slice. Often the honest one-month answer if licensing, data governance, and quality allow.

\textbf{The serving bill arrives regardless of method.} Whatever makes 64K \textit{work}, serving it costs:
\begin{itemize}
    \item \textbf{KV cache: $8\times$ per request.} For a 70B GQA model, per-token KV is $\approx 0.3$ MB, so a 64K request holds $\approx$20 GB of KV (vs.\ 2.5 GB at 8K)---concurrency per GPU collapses unless you add KV quantization, more HBM, or offloading. Do this arithmetic unprompted; the formula and worked example are in Chapter~\ref{chap:inference_optimization}.
    \item \textbf{Prefill compute grows $\sim$64$\times$} in the quadratic attention term: time-to-first-token moves from sub-second to seconds. Chunked prefill and prefix caching become mandatory, not optional.
    \item Long requests also occupy their KV for longer, compounding the throughput hit.
\end{itemize}

\textbf{Recommendation shape:} ship Option 1 behind a flag now, land Option 2 within the month as the supported path, put Option 3 on the roadmap, and route retrieval-style traffic to Option 4 so the expensive window is spent only on holistic tasks.

\redflags
\begin{itemize}
    \item ``Set \texttt{max\_position\_embeddings} to 65536 and redeploy''---positions past 8K are out-of-distribution rotary angles; attention beyond the trained range is garbage without one of the options above.
    \item Quoting perplexity as evidence that the extension worked.
    \item No mention of serving cost---the KV and prefill bill is often the binding constraint, not model quality.
\end{itemize}

\interviewtest{Do you know the extension-method menu and its quality ladder, that the native training recipe---not YaRN---is the frontier answer, and do you connect a context-length decision to its KV-cache and prefill consequences?}

\goodgreat
\begin{itemize}
    \item \badgeLfive{L5} Names PI/NTK/YaRN, knows they need brief fine-tuning, and picks YaRN for $8\times$.
    \item \badgeLsix{L6} Lays out the full option space including RAG and switching models, proposes an eval gate (NIAH as smoke test, RULER-style aggregation at 64K, short-context regression suite), and does the KV arithmetic.
    \item \badgeLseven{L7} Proposes the staged derisking plan (training-free now, YaRN this month, continued pretraining on the roadmap), quantifies the concurrency/TTFT impact on serving economics, and routes by workload---holistic tasks get the window, retrieval-style tasks get RAG.
\end{itemize}

\followups{``Why does NTK-aware scaling preserve short-range quality better than PI?'' ``How much long data does the continued-pretraining stage need?'' ``What breaks first if you serve 64K with the KV cache quantized to INT4?''}
\end{interviewq}

%--- Q2: Long-context quality degradation debugging ---

\begin{interviewq}
\textbf{Q: Your model's quality degrades beyond 32K context. Walk me through the diagnosis.}

\badgeLsix{L6} \quad \badgetype{Debugging} \quad \badgefreq{\freqmed}

\stronganswer

\textbf{Characterize the failure before touching anything.} Three discriminating questions: (1) Is the drop a \textit{cliff} at a specific length or a gradual slide? (2) Does it depend on \textit{where} in the context the needed information sits? (3) Does it reproduce with the serving stack taken out of the loop? The answers separate four root-cause families:

\textbf{Hypothesis 1: Position OOD (cliff at a fixed length, task-independent).} The model's trained-or-extended positional range effectively ends near 32K: a YaRN/NTK-extended model being served past its extension target, or---the classic production bug---the serving stack silently ignoring the \texttt{rope\_scaling} config or applying different scaling parameters than fine-tuning used. Signature: a NIAH heatmap goes uniformly red past a fixed length column, at every depth. Fix: audit the serving config against the training config first; only genuine extension (Section~\ref{subsec:inference_time_extension}) or retraining fixes a real range limit.

\textbf{Hypothesis 2: Lost in the middle (gradual, position-dependent).} Accuracy depends on answer position---the U-shaped curve, up to $\sim$20 points between mid-context and edge placement on multi-document QA (Section~\ref{subsec:long_context_eval}). Signature: permute the needle/answer position at a \textit{fixed} length; strong depth dependence with weak length dependence. Fix: reranker-aware ordering (most relevant first and last), fewer better chunks, or a model with a native long-context training stage.

\textbf{Hypothesis 3: Serving-side KV compression (gradual, worsens with length).} KV-cache quantization (INT8/INT4) accumulates error over long sequences, and KV eviction---sliding windows or token-dropping heuristics---silently discards distant tokens. A 32K sliding window produces \textit{exactly} ``fine to 32K, degrades beyond.'' Signature: rerun the same prompts offline with full-precision, full-KV inference; if the gap closes, the model is innocent. Check \texttt{kv\_cache\_dtype}, sliding-window settings, and any cache-eviction policy in the serving config.

\textbf{Hypothesis 4: Eval artifact.} The harness truncates inputs (tokenizer \texttt{max\_length}, template limits) so the needle never reaches the model; the chat template inflates token counts past the window; or the $>$32K test slice is drawn from a different, harder distribution (only certain domains produce 32K+ documents). Signature: log the exact token count entering the model and verify the needle survives preprocessing; compare task difficulty across length buckets.

\textbf{Diagnostic order---cheapest first:} (1) eval harness + serving/training config diff (minutes); (2) full-precision, full-KV ablation (hours); (3) position-permutation NIAH grid (hours); (4) only then model-side work---extension fine-tuning or a long-context training stage (days to weeks).

\redflags
\begin{itemize}
    \item Jumping straight to ``fine-tune on longer data''---the two most common real causes (config mismatch, serving-side KV policy) are fixable in an afternoon.
    \item Diagnosing from a single aggregate score without a length $\times$ depth breakdown.
    \item Not knowing what KV precision and eviction policy their own serving stack runs.
\end{itemize}

\interviewtest{Systematic isolation: model vs.\ serving stack vs.\ harness. Do you know the distinct signatures of position OOD, lost-in-the-middle, and KV compression, and do you order experiments by cost?}

\goodgreat
\begin{itemize}
    \item \badgeLfive{L5} Names lost-in-the-middle and out-of-distribution positions as candidate causes.
    \item \badgeLsix{L6} Adds the serving-side (KV quantization/eviction, rope-scaling config drift) and eval-artifact families, with a discriminating experiment for each.
    \item \badgeLseven{L7} Runs the decision tree in cost order, invokes the claimed-vs-effective context distinction, and institutionalizes the fix: a NIAH/RULER length-grid regression suite in CI so the next silent config change is caught before production.
\end{itemize}

\followups{``The NIAH grid is green but users still complain at 50K---what does that tell you?'' ``How would you detect a lost-in-the-middle regression in production without labels?'' ``Which KV quantization schemes degrade gracefully with length?''}
\end{interviewq}

%--- Q3: Long context vs RAG cost estimation ---

\begin{interviewq}
\textbf{Q: Estimate the cost per query for a RAG system vs a long-context LLM (128K).}

\badgeLsix{L6} \quad \badgetype{Estimation} \quad \badgefreq{\freqhigh}

\stronganswer

\textbf{Scenario:} Enterprise Q\&A over a knowledge base. User asks a question; the system returns an answer grounded in documents. Work with \textit{illustrative} prices and say so out loud: take $p_{\text{in}} = \$10$/M input tokens, $p_{\text{out}} = \$30$/M output tokens, and cached input at $\sim 0.1 \times p_{\text{in}} = \$1$/M as round numbers. API prices fall every year (most 2026 frontier prices are below these), so the durable content of this answer is the \textit{structure} of the comparison---input tokens dominate, caching rescales them---not the constants.

\textbf{Long-context approach:} Stuff the full context window with relevant documents.
\begin{itemize}
    \item Input: 128K tokens of context + $\sim$200 tokens of query/instructions = $\sim$128K input tokens.
    \item Output: $\sim$500 tokens (typical answer length).
    \item Cost at \$10/M input tokens, \$30/M output tokens: $(128{,}000 \times \$10 / 10^6) + (500 \times \$30 / 10^6) = \$1.28 + \$0.015 = \$1.30$ per query.
    \item Latency: 128K prefill is \textit{not} fast---prefill is compute-bound. For a 70B-class model, one forward pass over 128K tokens costs $\approx 2 \times 70 \times 10^9 \times 131{,}072 \approx 1.8 \times 10^{16}$ FLOPs; on 8$\times$H100 at 40\% MFU ($\approx 3.2 \times 10^{15}$ FLOPS) that is $\approx$5--6s of prefill, and real API time-to-first-token at 128K is commonly $\sim$10--30s. Add $\sim$500 tokens at 50 tok/sec $\approx$ 10s decode, for $\sim$15--40s total.
\end{itemize}

\textbf{RAG approach:} Retrieve top-$k$ chunks, stuff only those. (Both arms below are priced at an
illustrative \$10 per million input tokens and \$30 per million output tokens. List prices move every
few months and differ several-fold across model tiers---state your assumption, and note that the cost
\emph{ratio} is what the architecture decision turns on, not the absolute figures.)
\begin{itemize}
    \item Embedding the query: $\sim$1ms (cached model, single embedding).
    \item ANN retrieval: $\sim$5ms.
    \item Reranking (cross-encoder, top 50 $\rightarrow$ top 5): $\sim$50ms.
    \item Input to LLM: 5 chunks $\times$ 512 tokens $= 2{,}560$ tokens + 200 tokens query $= 2{,}760$ tokens.
    \item Output: $\sim$500 tokens.
    \item LLM cost: $(2{,}760 \times \$10/10^6) + (500 \times \$30/10^6) = \$0.028 + \$0.015 = \$0.043$ per query.
    \item Embedding + retrieval infra: $\sim$\$0.002 per query (amortized).
    \item \textbf{Total RAG cost: $\sim$\$0.045 per query.}
    \item Latency: retrieval ($\sim$55ms) + LLM prefill of 2,760 tokens ($\sim$100--300ms in practice) + decode ($\sim$10s) $\approx$ 10.2--10.4s---and time-to-first-token is well under a second.
\end{itemize}

\textbf{Comparison:}
\begin{itemize}
    \item \textbf{Cost ratio:} Long context is $\sim$29$\times$ more expensive (\$1.30 vs.\ \$0.045).
    \item \textbf{At 10K queries/day:} Long context = \$13,000/day. RAG = \$450/day. Difference: \$12,550/day = \$4.6M/year.
    \item \textbf{Latency:} Not comparable where users feel it: RAG's time-to-first-token is sub-second, while the long-context approach pays multi-second (often 10s+) prefill before the first token appears. Fast TTFT is a genuine RAG advantage; total wall-clock converges only for long, decode-dominated answers, and prefix caching is what rescues long context in multi-turn sessions.
    \item \textbf{Quality:} Long context is better for cross-document reasoning. RAG is better when the answer is localized in 1--2 passages.
\end{itemize}

\textbf{When long context is still worth the cost:} Single-document analysis (80-page contract), cross-document reasoning that RAG would miss, and low-volume high-value use cases where quality justifies cost.

\textbf{Caching changes the math---quantify it.} APIs price cached input separately, at roughly $0.1\times$ the fresh-input rate (some charge a small premium, $\sim$1.25$\times$, to \textit{write} the cache). For a multi-turn session over the same 128K document set: turn 1 pays the full $\sim$\$1.28 and writes the cache; every later turn pays $128{,}000 \times \$1/10^6 \approx \$0.13$ for the cached prefix (plus output)---a $10\times$ reduction, and cache hits also skip prefill compute, so TTFT collapses too. Amortized over a 10-turn session: $(\$1.28 + 9 \times \$0.13)/10 \approx \$0.25$ per turn---no longer $29\times$ RAG's \$0.045 but $\sim$5--6$\times$, close enough that quality (cross-document reasoning) rather than cost should decide. Any caching argument that does not use the cached-input price is hand-waving.

\redflags
\begin{itemize}
    \item Comparing only model inference cost without including RAG infrastructure (embedding, vector DB, reranking).
    \item Forgetting about caching---it dramatically changes the long-context cost for multi-turn use cases.
    \item Not converting to an annual dollar figure---the interviewer wants to see you reason about business impact.
\end{itemize}

\interviewtest{Can you do back-of-envelope cost estimation and convert it to business terms? Do you know the key cost drivers (input tokens dominate for long context, infrastructure for RAG)?}

\goodgreat
\begin{itemize}
    \item \badgeLfive{L5} Correctly estimates that long context is significantly more expensive per query and gives approximate numbers.
    \item \badgeLsix{L6} Provides detailed per-query cost breakdowns for both approaches, converts to annual cost at scale, and discusses caching as a cost modifier.
    \item \badgeLseven{L7} Factors in quality differences (long context wins for cross-document reasoning, RAG wins for localized answers), discusses the Pareto frontier of cost vs.\ quality, and proposes a hybrid architecture that routes queries to the cheaper option based on query type.
\end{itemize}

\followups{``How does caching change the cost for multi-turn conversations?'' ``What about latency comparison?'' ``When is the cost premium of long context justified?''}
\end{interviewq}

%--- Q4: The context-versus-retrieval decision ---

\begin{interviewq}
\textbf{Q: Long context window vs RAG---when do you choose each?}

\badgeLfive{L5} \quad \badgetype{Trade-off} \quad \badgefreq{\freqhigh}

\stronganswer

This is not an either/or decision. The right choice depends on four factors: corpus size, task type, cost constraints, and freshness requirements.

\textbf{Choose long context when:}
\begin{itemize}
    \item The corpus fits in the window (single document, $<$100K tokens). Analyzing an 80-page legal contract at $\sim$40K tokens is a perfect long-context use case.
    \item The task requires holistic reasoning across the entire corpus---summarization, theme extraction, contradiction detection. RAG would miss connections between non-adjacent passages.
    \item Offline/batch analysis where cost per query is amortized and latency is non-critical.
\end{itemize}

\textbf{Choose RAG when:}
\begin{itemize}
    \item The corpus is large (10K+ documents). No context window can hold it all.
    \item Cost matters. At an illustrative \$10/M input tokens (state API prices as variables---they fall every year), 128K input = \$1.28/query vs.\ \$0.05/query for RAG with 5K tokens. At 10K queries/day: \$12,800/day vs.\ \$500/day. Cached-input pricing ($\sim$0.1$\times$ the fresh rate) shrinks the repeated-prefix case by $\sim$10$\times$ but does not close the gap for independent queries.
    \item The knowledge base updates frequently. Updating a vector index is cheap.
    \item You need source attribution. RAG provides ``answer from Document X, page Y'' naturally.
\end{itemize}

\textbf{Combine both (the strongest production pattern):}
\begin{itemize}
    \item RAG retrieves top 20--30 chunks ($\sim$10K--15K tokens). A long-context model reasons across them jointly. You get retrieval precision + cross-chunk reasoning without stuffing 128K tokens.
    \item For multi-hop questions: retrieve, generate intermediate answer, retrieve again, generate final answer (``retrieve-and-reason'' loop).
\end{itemize}

\textbf{The ``lost in the middle'' problem:} LLMs pay most attention to the beginning and end of context---up to $\sim$20 points of multi-document QA accuracy separate mid-context from edge placement (partially fixed by modern long-context training; Section~\ref{subsec:long_context_eval}). For long context: order documents with the most relevant first and last. For RAG: the reranker's ordering directly impacts generation quality.

\redflags
\begin{itemize}
    \item ``Just use the longest context window available''---cost and lost-in-the-middle make this impractical at scale.
    \item Not discussing cost---it is the deciding factor for most enterprise decisions.
    \item ``RAG is always better''---long context wins for holistic cross-document reasoning and single-document analysis.
\end{itemize}

\interviewtest{Nuanced engineering judgment. Can you reason about cost, latency, and quality simultaneously? Do you know the lost-in-the-middle problem? Do you recognize the hybrid approach?}

\goodgreat
\begin{itemize}
    \item \badgeLfive{L5} Correctly identifies corpus size and cost as the primary decision factors.
    \item \badgeLsix{L6} Provides concrete cost estimates, discusses lost-in-the-middle, and recommends the hybrid approach.
    \item \badgeLseven{L7} Factors in caching (prefix caching changes the cost equation for multi-turn), reasons about the latency difference ($O(n^2)$ attention for long context vs.\ fast prefill for RAG), and proposes an intelligent routing system that sends queries to long context or RAG based on query characteristics.
\end{itemize}

\followups{``What is the lost-in-the-middle problem?'' ``How does caching change the cost comparison?'' ``What about latency at 128K tokens?''}
\end{interviewq}

\begin{keyinsight}[Chapter Summary: A Window Is Capacity, Not a Guarantee]
Two things make this chapter's material interview-ready. \textbf{How a window is built and verified}: the 2026 answer to ``give this model 128K'' is a continued-pretraining recipe---raise the RoPE base, extend length in stages on a long-data mixture, finish with long-context SFT---with position interpolation, NTK-aware scaling, and YaRN as the retrofits you reach for when you cannot retrain, and with any claimed length treated as marketing until your own RULER-style length$\times$depth grid says otherwise. \textbf{When to spend the window instead of retrieving}: size the corpus against the window, check the effective length, then decide on task shape and failure mode---position-and-length degradation on one side, retrieval recall on the other---with token cost and the standing cost of owning a pipeline as the tiebreak, and the hybrid (retrieve wide, rerank hard, reason over a filled 10K--30K window) as the production default. The pipeline that hybrid runs on is [NLP~9]; the retrieval infrastructure under it is [SR~3], [SR~4], and [SR~8]. The three ways candidates lose this chapter: quoting YaRN as the state of the art, accepting a claimed context length, and comparing a prompt to a system on token price alone.
\end{keyinsight}
